\kbvcontinued ಪದ. ಇದು “ಕತ್ತರಿ”ಯಾಗಿ ಮಾರ್ಪಾಡು ಹೊಂದಿ ಕನ್ನಡದಲ್ಲಿ ಬಳಕೆಗೆ ಬಂದಿದೆ ; ಆದ್ದರಿಂದ ಇದನ್ನು ಕರ್ತರಿ ಶಬ್ದದ ತದ್ಭವ ರೂಪ ಎನ್ನುತ್ತೇವೆ. ಹೆಚ್ಚಿನ ಉದಾಹರಣೆಗಳನ್ನು ನೋಡೋಣ.\par

\vspace{5pt plus 4pt minus 1pt}
\begin{center}
{\renewcommand{\arraystretch}{1.22}\setlength{\tabcolsep}{5.5pt}
\begin{tabularx}{\textwidth}{X X | X X | X X}
{\kbvsecfont ತತ್ಸಮ} & {\kbvsecfont ತದ್ಭವ} & {\kbvsecfont ತತ್ಸಮ} & {\kbvsecfont ತದ್ಭವ} & {\kbvsecfont ತತ್ಸಮ} & {\kbvsecfont ತದ್ಭವ} \\[3pt]
ಆಶಾ & ಆಸೆ & ಗಂಗಾ & ಗಂಗೆ & ರಾಜನ್‌ & ರಾಜ \\
ಗ್ರಂಥಿ & ಗಂಟು & ಯೋಗಿ & ಜೋಗಿ & ಕುಠಾರ & ಕೊಡಲಿ \\
ಚಂದ್ರ & ಚಂದಿರ & ಶೂನ್ಯ & ಸೊನ್ನೆ & ಶರ್ಕರ & ಸಕ್ಕರೆ \\
ಧೃತಿ & ದಠಿ & ಸಂಸ್ಥಾ & ಸಂಸ್ಥೆ & ಕಾಮಿನೀ & ಕಾಮಿನಿ \\
ಬ್ರಹ್ಮ & ಬೊಮ್ಮ & ಮುಕ್ತಿ & ಮುಕುತಿ & ತೃತೀಯಾ & ತದಿಗೆ \\
ಮಾಲಾ & ಮಾಲೆ & ಘಂಟಾ & ಗಂಟೆ & ಯಶಸ್ & ಯಶಸ್ಸು \\
ಯಮ & ಜವ & ವಸತಿ & ಬಸದಿ & ರಾಕ್ಷಸ & ರಕ್ಕಸ \\
ವೀಥಿ & ಬೀದಿ & ಚತುರ & ಚದುರ & ವಿಸ್ತಾರ & ಬಿತ್ತರ \\
\end{tabularx}}
\end{center}
\vspace{5pt plus 4pt minus 1pt}

\kbvexercise{ಅಭ್ಯಾಸ ಪತ್ರಿಕೆ ೧೮}

\vspace{2.5pt plus 3pt minus 0.5pt}\noindent\hspace*{1.2em}\kbvnum{(1)}\hspace{0.55em}ತತ್ಸಮ, ತದ್ಭವ—ಈ ಪದಗಳ ಅರ್ಥವನ್ನು ವಿವರಿಸಿ.\par

\vspace{2.5pt plus 3pt minus 0.5pt}\noindent\hspace*{1.2em}\kbvnum{(2)}\hspace{0.55em}ಕೆಳಗೆ ಕಾಣಿಸಿರುವ ಪದಗಳ ತತ್ಸಮ ರೂಪಗಳನ್ನು ತಿಳಿಸಿ : ಕನ್ನಿಕೆ, ಸಿಂಗರ, ಪಯಣ, ನೇಹ, ಸಿಂಗ, ಲಕ್ಷಿ, ಸೀತೆ, ಲತೆ.\par

\vspace{2.5pt plus 3pt minus 0.5pt}\noindent\hspace*{1.2em}\kbvnum{(3)}\hspace{0.55em}ಕೆಳಗೆ ಕಾಣಿಸಿರುವ ಪದಗಳ ತದ್ಭವ ರೂಪಗಳನ್ನು ತಿಳಿಸಿ : ವಿದ್ವಾಂಸ, ಚತುರ್ಥೀ, ಕುಮಾರೀ, ಶಿರಸ್, ಸುಧಾ, ವಸುಂಧರಾ.\par

\vspace{2.5pt plus 3pt minus 0.5pt}\noindent\hspace*{1.2em}\kbvnum{(4)}\hspace{0.55em}ಕನ್ನಡದಲ್ಲಿ ಬಳಕೆಯಲ್ಲಿರುವ ಹತ್ತು ಅನ್ಯದೇಶೀಯ ಪದಗಳನ್ನು ಹೆಸರಿಸಿ. ಉದಾ : ದವಾಖಾನೆ, ಪಾದ್ರಿ.... ಇತ್ಯಾದಿ.\par

\kbvchapter{ಪಾಠ ೧೯}{ಲೇಖನ ಚಿಹ್ನೆಗಳು}

\vspace{5pt plus 4pt minus 1pt}
\noindent{\renewcommand{\arraystretch}{1.24}\setlength{\tabcolsep}{4pt}
\begin{tabularx}{\textwidth}{X X X X}
{\kbvsecfont\color{kbvaccentdark}ಪೂರ್ಣ ವಿರಾಮ ಚಿಹ್ನೆ} & {\kbvsecfont\color{kbvaccentdark}{[}.{]}} & {\kbvsecfont\color{kbvaccentdark}ಭಾವಸೂಚಕ ಚಿಹ್ನೆ} & {\kbvsecfont\color{kbvaccentdark}{[}!{]}} \\[1pt]
\arrayrulecolor{kbvline}\specialrule{0.5pt}{0pt}{2pt}\arrayrulecolor{black}
ಅರ್ಧ ವಿರಾಮ ಚಿಹ್ನೆ & {[};{]} & ಉದ್ಧರಣ ಚಿಹ್ನೆ & {[}“ ”{]} \\
\end{tabularx}}
\vspace{5pt plus 4pt minus 1pt}
